\section{Вибір криптографічного механізму}

\subsection{Ситуація 1. Захист електронного документа}

Умова: Необхідно забезпечити можливість перевірити, 
що електронний документ був створений певним власником закритого 
ключа та не був модифікований після підписання.

Запропонований криптографічний механізм: Цифровий підпис та хешування.

Можливий алгоритм або схема: Ed25519.

Обґрунтування вибору: Спочатку для документа обчислюється хеш, 
після чого він підписується приватним ключем власника. Для перевірки 
використовується відповідний відкритий ключ. Якщо документ було змінено 
після підписання, перевірка підпису завершиться невдало. Ed25519 забезпечує 
високу швидкість роботи та компактний розмір підпису. При цьому для підтвердження 
особи власника необхідно також мати надійний зв'язок між його особою та відкритим ключем.

\subsection{Ситуація 2. Формування спільного секретного ключа}

Умова: Alice та Bob повинні сформувати спільний секретний ключ, обмінюючись 
інформацією через недовірену мережу.

Запропонований криптографічний механізм: Узгодження ключів.

Можливий алгоритм або схема: ECDH на основі еліптичної кривої X25519.

Обґрунтування вибору: Кожна сторона генерує власну пару ключів і обмінюється відкритими ключами. 
Після цього Alice та Bob незалежно обчислюють однаковий спільний секрет, використовуючи власний приватний 
ключ та відкритий ключ іншої сторони. Сам секрет через мережу не передається. X25519 забезпечує ефективне 
узгодження ключів, однак для захисту від атак типу <<людина посередині>> необхідна автентифікація сторін.

\subsection{Ситуація 3. Конфіденційна передача великого файла}

Умова: Необхідно передати файл розміром 2 ГБ у конфіденційному вигляді. Відкритий ключ отримувача відомий відправнику.

Запропонований криптографічний механізм: Гібридне шифрування, яке поєднує асиметричну криптографію та симетричне шифрування.

Можливий алгоритм або схема: Для встановлення секретного ключа можна використати RSA-OAEP, а для шифрування файла — AES-256-GCM.

Обґрунтування вибору та переваги гібридної схеми: Спочатку відправник генерує випадковий симетричний ключ, 
яким шифрує файл за допомогою AES-256-GCM. Сам симетричний ключ шифрується відкритим ключем отримувача за допомогою RSA-OAEP. 
Отримувач розшифровує ключ своїм приватним ключем, а потім використовує його для відновлення початкового файла.
Безпосереднє асиметричне шифрування всього файла не є оптимальним, оскільки такі алгоритми мають обмеження на розмір даних, 
які можна зашифрувати за одну операцію, та працюють повільніше за симетричні. Гібридна схема дозволяє швидко зашифрувати великий 
обсяг інформації, використовуючи асиметричний алгоритм лише для захисту невеликого симетричного ключа.

\subsection{Ситуація 4. Безпечне зберігання паролів}

Умова: Сервер повинен безпечно зберігати паролі користувачів та перевіряти їх під час автентифікації.

Чи потрібна асиметрична криптографія: Ні, для цього завдання асиметрична криптографія не потрібна. Паролі не потрібно 
шифрувати для подальшого відновлення, достатньо зберігати їхні спеціально обчислені хеші.

Запропонований механізм: Хешування паролів із використанням індивідуальної випадкової солі.

Обґрунтування вибору: Можна використати алгоритм Argon2id, який спеціально призначений для захищеного зберігання паролів. 
Під час реєстрації для кожного пароля генерується випадкова сіль, після чого обчислюється його хеш. На сервері зберігаються сіль 
та отриманий хеш, але не сам пароль. Під час автентифікації сервер повторно обчислює хеш введеного пароля з тією самою сіллю 
та порівнює результати. Argon2id ускладнює масовий підбір паролів завдяки використанню обчислювальних ресурсів і пам'яті.

\subsection{Ситуація 5. Автентичність оновлень програмного забезпечення}

Умова: Виробник випускає оновлення програмного забезпечення для великої кількості пристроїв. Створювати справжні оновлення може 
лише виробник, а кожен пристрій повинен мати можливість перевірити їх автентичність.

Запропонований криптографічний механізм: Цифровий підпис та хешування.

Можливий алгоритм або схема: Ed25519.

Розподіл ключів між виробником та пристроями: Виробник зберігає приватний ключ у захищеному середовищі та використовує 
його для підписання оновлень. На кожному пристрої зберігається відповідний відкритий ключ, який використовується для перевірки підпису. 
Приватний ключ виробника ніколи не передається пристроям.

Обґрунтування вибору: Перед розповсюдженням оновлення виробник створює його цифровий підпис. Пристрій після отримання оновлення 
перевіряє підпис за допомогою відкритого ключа. Якщо перевірка успішна, це підтверджує, що оновлення було підписане власником відповідного 
приватного ключа та не було змінене після підписання. Такий механізм дозволяє поширювати оновлення серед великої кількості пристроїв без передавання їм секретного ключа.
